% 猎户座（综述）
% license CCBYSA3
% type Wiki

本文根据 CC-BY-SA 协议转载翻译自维基百科\href{https://en.wikipedia.org/wiki/Orion_(constellation)}{相关文章}。

猎户座是在北天球冬季可见的一组显著恒星。它是现代划分的88个星座之一，同时也名列公元2世纪天文学家托勒密所列出的48个星座之中。其名称源自希腊神话中的一位猎人。

在北半球的冬季夜晚，猎户座最为醒目，另外还有五个星座与其一同构成“冬季六边形”星群。猎户座中最明亮的两颗恒星是参宿七（$\beta$）和参宿四（$\alpha$），二者均位列夜空中最亮的恒星之列；它们都是超巨星，并且亮度略有变化。此外，猎户座中还有六颗视星等亮于3.0的恒星，其中三颗构成了著名的猎户座腰带这一短而笔直的星群。猎户座还是一年一度猎户座流星雨的辐射点，该流星雨是与哈雷彗星相关的最强流星雨之一；同时，猎户座内还包含猎户座星云，这是夜空中最明亮的星云之一。
\subsection{特征}
\begin{figure}[ht]
\centering
\includegraphics[width=6cm]{./figures/c950af7d2bff8c26.png}
\caption{猎户座的星群示意图，如肉眼所见，并绘制了连线。} \label{fig_Orion_1}
\end{figure}
猎户座的西北方毗邻金牛座，西南方为波江座，南方为天兔座，东方为麒麟座，东北方为双子座。猎户座覆盖约594平方度，在88个星座中面积排名第26位。该星座的边界由比利时天文学家欧仁·德尔波特于1930年划定，其边界由一个26边形构成。在赤道坐标系中，这些边界的赤经范围介于 $04^{h}43.3^{m}$与$06^{h}25.5^{m}$之间，而赤纬范围则介于 22.87° 与 −10.97° 之间。$^{[2]}$ 国际天文学联合会于1922年采用的该星座三字母缩写为 “Ori”。$^{[3]}$

猎户座在一月至四月的夜空中最为清晰可见，$^{[4]}$ 此时北半球正值冬季，而南半球为夏季。在热带地区（距赤道约8°以内），猎户座会通过天顶。

从五月到七月（北半球的夏季、南半球的冬季），猎户座位于白昼天空中，因此在大多数纬度地区不可见。然而，在南半球冬季的数月中，南极洲大部分地区即使在正午太阳也位于地平线以下。此时，恒星（以及猎户座，但仅限于最亮的恒星）会在地方正午前后数小时内于暮光中可见，出现在北方低空、天空中最明亮的区域，太阳正位于地平线下方不远处。在同一时刻、南极点本身（阿蒙森—斯科特南极站），参宿七仅高出地平线8°，而猎户座腰带几乎沿着地平线掠过。南半球的夏季月份中，猎户座在其他地区通常可于夜空中观测，但在南极洲却完全不可见，因为在南极圈以南的这一时期太阳并不落下。$^{[5][6]}$

在接近赤道的国家（例如肯尼亚、印度尼西亚、哥伦比亚、厄瓜多尔），猎户座在12月午夜前后位于头顶上方，而在2月则出现在傍晚的天空中。
\subsection{导航用途}
\begin{figure}[ht]
\centering
\includegraphics[width=6cm]{./figures/96807e4588169a15.png}
\caption{利用猎户座定位邻近星座中的恒星} \label{fig_Orion_2}
\end{figure}
猎户座在定位其他恒星时极为有用。将猎户座腰带的连线向东南方向延伸，可以找到天狼星（$\alpha$ CMa）；向西北方向延伸，则可找到毕宿五（$\alpha$ Tau）。沿着猎户座两肩向东方画一条直线，可指示南河三（$\alpha$ CMi）的方向。从参宿七穿过参宿四所作的连线则指向双子座的北河二与北河三（$\alpha$ Gem 与 $\beta$ Gem）。此外，参宿七还是“冬季圆环”星群的一部分。天狼星和南河三也可以通过从猎户座出发的假想连线加以定位（见星图），它们同时也是“冬季三角形”和“冬季圆环”的组成恒星。$^{[7]}$
\subsection{特征}
\begin{figure}[ht]
\centering
\includegraphics[width=6cm]{./figures/fa98c4485118af07.png}
\caption{约1825年在伦敦出版的一套星图卡片《乌拉尼娅之镜》中所描绘的猎户座} \label{fig_Orion_3}
\end{figure}
猎户座中最明亮的七颗恒星在夜空中构成了一个极具辨识度的沙漏形星群或图案。其中四颗恒星——参宿七（Rigel）、参宿四（Betelgeuse）、参宿五（Bellatrix）和参宿六（Saiph）——形成一个近似矩形的大轮廓，而在其中心位置则是由参宿一（Alnitak）、参宿二（Alnilam）和参宿三（Mintaka）组成的猎户座腰带三星。猎户的头部由另一颗第八颗恒星标示，即觜宿一（Meissa），对观测者而言也相当明亮。自腰带向下延伸的是一条由三颗天体构成的较短直线，即猎户座之剑（其中间的“星体”实际上并非恒星，而是猎户座星云），亦称为猎人的佩剑。

其中许多恒星都是高光度、炽热的蓝色超巨星，腰带与佩剑中的恒星共同构成了猎户座 OB1 协会。参宿四以其显著的红色光芒格外醒目，但它可能仍是同一恒星集团中的一颗逃逸成员。

按距离排列的猎户座恒星（红—绿三维视图）及各恒星的亮度（以恒星大小表示）  
\subsubsection{亮星}
